% Drop-in section for the E2 corpus measurement.
% To include: put % Drop-in section for the E2 corpus measurement.
% To include: put % Drop-in section for the E2 corpus measurement.
% To include: put % Drop-in section for the E2 corpus measurement.
% To include: put \input{section_e2} immediately before \section{Mechanistic fingerprint...}
% Costs roughly half a page. Frees space by trimming Related Work (Sec. 6) and the
% second paragraph of Sec. 3's "Behavioral reliance" discussion.

\section{What the corpora make predictable}
\label{sec:corpus}

The account above rests on \emph{predictive necessity}: because the pretraining objective is
masked-cell reconstruction, a corpus can only teach cross-table computation if its masked cells
are hard to recover from their own row. That is a property of the data, not of any model, so we
measure it directly and without involving the RT at all.

For each corpus we sample 40 databases, recover the FK graph, and for each target column fit
identical gradient-boosted regressors three times: from the target's own row; from its own row
plus the columns of its FK parents; and from its own row plus aggregates over its FK children.
The gain from adding relational features is the neighbour predictive information in that corpus.
Splitting parents from children matters because the two map onto different routes in the RT: a
cell reading its parents flows through the \texttt{feat} stream's parent-cell inclusion, while
aggregation over children is what the \texttt{nbr} stream does. Scores are clipped to $[-1,1]$
before differencing, since unbounded negative $R^2$ lets one diverged fit dominate a mean, and we
report medians with 95\% bootstrap intervals.

\begin{table}[t]
\caption{Neighbour predictive information per corpus, measured with no RT involved (40 databases
each; median, 95\% bootstrap CI). RelDiff is the only corpus in which reading a neighbouring table
measurably helps recover a held-out cell. RDB-PFN is its mirror image: the most learnable
within-row structure of any corpus and no relational signal at all.}
\label{tab:e2}
\centering
\small
\begin{tabular}{lcccc}
\toprule
Corpus & Tasks & Own-row $R^2$ & Parent lift & Child lift \\
\midrule
RelDiff & 48  & 0.022 & \textbf{+0.072 [+0.018, +0.091]} & $-$0.001 \\
PluRel  & 113 & 0.166 & +0.009 [+0.001, +0.022] & $-$0.000 \\
RDB-PFN & 120 & 0.436 & +0.001 [$-$0.004, +0.009] & $-$0.005 \\
GRDM    & 120 & $-$0.011 & $-$0.057 [$-$0.076, $-$0.044] & $-$0.105 \\
\bottomrule
\end{tabular}
\end{table}

Table~\ref{tab:e2} orders the generators the same way every downstream measurement does, and it
does so before any model is trained. RelDiff's parent lift is eight times PluRel's and roughly
seventy times RDB-PFN's, whose interval spans zero. Read as a share of the predictable signal,
neighbours contribute 77\% for RelDiff, 5\% for PluRel and 0.1\% for RDB-PFN. RelDiff's cells are
close to unrecoverable from their own row ($R^2$ 0.022) and become recoverable once a parent is
visible, which is precisely the condition under which masked-cell reconstruction rewards learning
to read across tables. RDB-PFN's profile --- the most predictable rows and no relational signal
--- is the data-side counterpart of its benchmark behaviour, strong at short context and untouched
by ablation.

Two qualifications. Child lift is flat in every corpus, RelDiff included, so what its data
rewards specifically is reading parents; the transfer to child-aggregation tasks at benchmark time
is not something the corpus directly incentivises, and our aggregates (mean, standard deviation,
count) may simply be too coarse to capture it. Second, GRDM's negative lift is best read alongside
a property of its shipped corpus rather than its design: every string and datetime column was
converted to a float and never converted back, so its databases carry no text or categorical
columns at all (0 against RelDiff's 234 over the same references). The RT's text encoder and
MiniLM pathway therefore received no signal in that arm, which is a confound worth stating
plainly: the GRDM arm differs from the others in more than its generator.

\paragraph{Diversity is not the mechanism.}
A natural alternative is that RelDiff simply produces more diverse data. Measuring feature, table
and cross-table diversity directly from the generated databases --- never through any of the four
models, so the measurement cannot be an artifact of what it explains --- RelDiff is indeed most
diverse on all three axes (1.10 / 0.57 / 1.04). But the ranking of the others does not track
mechanism formation: PluRel exceeds GRDM on cross-table diversity (0.85 vs.\ 0.71) yet learns no
usable relational computation (Table~\ref{tab:e7}), while GRDM does on in-distribution schemas.
Diversity describes a generator's marginal output; predictive necessity describes whether values
are \emph{coupled} to structure, and only the latter separates the four models.
 immediately before \section{Mechanistic fingerprint...}
% Costs roughly half a page. Frees space by trimming Related Work (Sec. 6) and the
% second paragraph of Sec. 3's "Behavioral reliance" discussion.

\section{What the corpora make predictable}
\label{sec:corpus}

The account above rests on \emph{predictive necessity}: because the pretraining objective is
masked-cell reconstruction, a corpus can only teach cross-table computation if its masked cells
are hard to recover from their own row. That is a property of the data, not of any model, so we
measure it directly and without involving the RT at all.

For each corpus we sample 40 databases, recover the FK graph, and for each target column fit
identical gradient-boosted regressors three times: from the target's own row; from its own row
plus the columns of its FK parents; and from its own row plus aggregates over its FK children.
The gain from adding relational features is the neighbour predictive information in that corpus.
Splitting parents from children matters because the two map onto different routes in the RT: a
cell reading its parents flows through the \texttt{feat} stream's parent-cell inclusion, while
aggregation over children is what the \texttt{nbr} stream does. Scores are clipped to $[-1,1]$
before differencing, since unbounded negative $R^2$ lets one diverged fit dominate a mean, and we
report medians with 95\% bootstrap intervals.

\begin{table}[t]
\caption{Neighbour predictive information per corpus, measured with no RT involved (40 databases
each; median, 95\% bootstrap CI). RelDiff is the only corpus in which reading a neighbouring table
measurably helps recover a held-out cell. RDB-PFN is its mirror image: the most learnable
within-row structure of any corpus and no relational signal at all.}
\label{tab:e2}
\centering
\small
\begin{tabular}{lcccc}
\toprule
Corpus & Tasks & Own-row $R^2$ & Parent lift & Child lift \\
\midrule
RelDiff & 48  & 0.022 & \textbf{+0.072 [+0.018, +0.091]} & $-$0.001 \\
PluRel  & 113 & 0.166 & +0.009 [+0.001, +0.022] & $-$0.000 \\
RDB-PFN & 120 & 0.436 & +0.001 [$-$0.004, +0.009] & $-$0.005 \\
GRDM    & 120 & $-$0.011 & $-$0.057 [$-$0.076, $-$0.044] & $-$0.105 \\
\bottomrule
\end{tabular}
\end{table}

Table~\ref{tab:e2} orders the generators the same way every downstream measurement does, and it
does so before any model is trained. RelDiff's parent lift is eight times PluRel's and roughly
seventy times RDB-PFN's, whose interval spans zero. Read as a share of the predictable signal,
neighbours contribute 77\% for RelDiff, 5\% for PluRel and 0.1\% for RDB-PFN. RelDiff's cells are
close to unrecoverable from their own row ($R^2$ 0.022) and become recoverable once a parent is
visible, which is precisely the condition under which masked-cell reconstruction rewards learning
to read across tables. RDB-PFN's profile --- the most predictable rows and no relational signal
--- is the data-side counterpart of its benchmark behaviour, strong at short context and untouched
by ablation.

Two qualifications. Child lift is flat in every corpus, RelDiff included, so what its data
rewards specifically is reading parents; the transfer to child-aggregation tasks at benchmark time
is not something the corpus directly incentivises, and our aggregates (mean, standard deviation,
count) may simply be too coarse to capture it. Second, GRDM's negative lift is best read alongside
a property of its shipped corpus rather than its design: every string and datetime column was
converted to a float and never converted back, so its databases carry no text or categorical
columns at all (0 against RelDiff's 234 over the same references). The RT's text encoder and
MiniLM pathway therefore received no signal in that arm, which is a confound worth stating
plainly: the GRDM arm differs from the others in more than its generator.

\paragraph{Diversity is not the mechanism.}
A natural alternative is that RelDiff simply produces more diverse data. Measuring feature, table
and cross-table diversity directly from the generated databases --- never through any of the four
models, so the measurement cannot be an artifact of what it explains --- RelDiff is indeed most
diverse on all three axes (1.10 / 0.57 / 1.04). But the ranking of the others does not track
mechanism formation: PluRel exceeds GRDM on cross-table diversity (0.85 vs.\ 0.71) yet learns no
usable relational computation (Table~\ref{tab:e7}), while GRDM does on in-distribution schemas.
Diversity describes a generator's marginal output; predictive necessity describes whether values
are \emph{coupled} to structure, and only the latter separates the four models.
 immediately before \section{Mechanistic fingerprint...}
% Costs roughly half a page. Frees space by trimming Related Work (Sec. 6) and the
% second paragraph of Sec. 3's "Behavioral reliance" discussion.

\section{What the corpora make predictable}
\label{sec:corpus}

The account above rests on \emph{predictive necessity}: because the pretraining objective is
masked-cell reconstruction, a corpus can only teach cross-table computation if its masked cells
are hard to recover from their own row. That is a property of the data, not of any model, so we
measure it directly and without involving the RT at all.

For each corpus we sample 40 databases, recover the FK graph, and for each target column fit
identical gradient-boosted regressors three times: from the target's own row; from its own row
plus the columns of its FK parents; and from its own row plus aggregates over its FK children.
The gain from adding relational features is the neighbour predictive information in that corpus.
Splitting parents from children matters because the two map onto different routes in the RT: a
cell reading its parents flows through the \texttt{feat} stream's parent-cell inclusion, while
aggregation over children is what the \texttt{nbr} stream does. Scores are clipped to $[-1,1]$
before differencing, since unbounded negative $R^2$ lets one diverged fit dominate a mean, and we
report medians with 95\% bootstrap intervals.

\begin{table}[t]
\caption{Neighbour predictive information per corpus, measured with no RT involved (40 databases
each; median, 95\% bootstrap CI). RelDiff is the only corpus in which reading a neighbouring table
measurably helps recover a held-out cell. RDB-PFN is its mirror image: the most learnable
within-row structure of any corpus and no relational signal at all.}
\label{tab:e2}
\centering
\small
\begin{tabular}{lcccc}
\toprule
Corpus & Tasks & Own-row $R^2$ & Parent lift & Child lift \\
\midrule
RelDiff & 48  & 0.022 & \textbf{+0.072 [+0.018, +0.091]} & $-$0.001 \\
PluRel  & 113 & 0.166 & +0.009 [+0.001, +0.022] & $-$0.000 \\
RDB-PFN & 120 & 0.436 & +0.001 [$-$0.004, +0.009] & $-$0.005 \\
GRDM    & 120 & $-$0.011 & $-$0.057 [$-$0.076, $-$0.044] & $-$0.105 \\
\bottomrule
\end{tabular}
\end{table}

Table~\ref{tab:e2} orders the generators the same way every downstream measurement does, and it
does so before any model is trained. RelDiff's parent lift is eight times PluRel's and roughly
seventy times RDB-PFN's, whose interval spans zero. Read as a share of the predictable signal,
neighbours contribute 77\% for RelDiff, 5\% for PluRel and 0.1\% for RDB-PFN. RelDiff's cells are
close to unrecoverable from their own row ($R^2$ 0.022) and become recoverable once a parent is
visible, which is precisely the condition under which masked-cell reconstruction rewards learning
to read across tables. RDB-PFN's profile --- the most predictable rows and no relational signal
--- is the data-side counterpart of its benchmark behaviour, strong at short context and untouched
by ablation.

Two qualifications. Child lift is flat in every corpus, RelDiff included, so what its data
rewards specifically is reading parents; the transfer to child-aggregation tasks at benchmark time
is not something the corpus directly incentivises, and our aggregates (mean, standard deviation,
count) may simply be too coarse to capture it. Second, GRDM's negative lift is best read alongside
a property of its shipped corpus rather than its design: every string and datetime column was
converted to a float and never converted back, so its databases carry no text or categorical
columns at all (0 against RelDiff's 234 over the same references). The RT's text encoder and
MiniLM pathway therefore received no signal in that arm, which is a confound worth stating
plainly: the GRDM arm differs from the others in more than its generator.

\paragraph{Diversity is not the mechanism.}
A natural alternative is that RelDiff simply produces more diverse data. Measuring feature, table
and cross-table diversity directly from the generated databases --- never through any of the four
models, so the measurement cannot be an artifact of what it explains --- RelDiff is indeed most
diverse on all three axes (1.10 / 0.57 / 1.04). But the ranking of the others does not track
mechanism formation: PluRel exceeds GRDM on cross-table diversity (0.85 vs.\ 0.71) yet learns no
usable relational computation (Table~\ref{tab:e7}), while GRDM does on in-distribution schemas.
Diversity describes a generator's marginal output; predictive necessity describes whether values
are \emph{coupled} to structure, and only the latter separates the four models.
 immediately before \section{Mechanistic fingerprint...}
% Costs roughly half a page. Frees space by trimming Related Work (Sec. 6) and the
% second paragraph of Sec. 3's "Behavioral reliance" discussion.

\section{What the corpora make predictable}
\label{sec:corpus}

The account above rests on \emph{predictive necessity}: because the pretraining objective is
masked-cell reconstruction, a corpus can only teach cross-table computation if its masked cells
are hard to recover from their own row. That is a property of the data, not of any model, so we
measure it directly and without involving the RT at all.

For each corpus we sample 40 databases, recover the FK graph, and for each target column fit
identical gradient-boosted regressors three times: from the target's own row; from its own row
plus the columns of its FK parents; and from its own row plus aggregates over its FK children.
The gain from adding relational features is the neighbour predictive information in that corpus.
Splitting parents from children matters because the two map onto different routes in the RT: a
cell reading its parents flows through the \texttt{feat} stream's parent-cell inclusion, while
aggregation over children is what the \texttt{nbr} stream does. Scores are clipped to $[-1,1]$
before differencing, since unbounded negative $R^2$ lets one diverged fit dominate a mean, and we
report medians with 95\% bootstrap intervals.

\begin{table}[t]
\caption{Neighbour predictive information per corpus, measured with no RT involved (40 databases
each; median, 95\% bootstrap CI). RelDiff is the only corpus in which reading a neighbouring table
measurably helps recover a held-out cell. RDB-PFN is its mirror image: the most learnable
within-row structure of any corpus and no relational signal at all.}
\label{tab:e2}
\centering
\small
\begin{tabular}{lcccc}
\toprule
Corpus & Tasks & Own-row $R^2$ & Parent lift & Child lift \\
\midrule
RelDiff & 48  & 0.022 & \textbf{+0.072 [+0.018, +0.091]} & $-$0.001 \\
PluRel  & 113 & 0.166 & +0.009 [+0.001, +0.022] & $-$0.000 \\
RDB-PFN & 120 & 0.436 & +0.001 [$-$0.004, +0.009] & $-$0.005 \\
GRDM    & 120 & $-$0.011 & $-$0.057 [$-$0.076, $-$0.044] & $-$0.105 \\
\bottomrule
\end{tabular}
\end{table}

Table~\ref{tab:e2} orders the generators the same way every downstream measurement does, and it
does so before any model is trained. RelDiff's parent lift is eight times PluRel's and roughly
seventy times RDB-PFN's, whose interval spans zero. Read as a share of the predictable signal,
neighbours contribute 77\% for RelDiff, 5\% for PluRel and 0.1\% for RDB-PFN. RelDiff's cells are
close to unrecoverable from their own row ($R^2$ 0.022) and become recoverable once a parent is
visible, which is precisely the condition under which masked-cell reconstruction rewards learning
to read across tables. RDB-PFN's profile --- the most predictable rows and no relational signal
--- is the data-side counterpart of its benchmark behaviour, strong at short context and untouched
by ablation.

Two qualifications. Child lift is flat in every corpus, RelDiff included, so what its data
rewards specifically is reading parents; the transfer to child-aggregation tasks at benchmark time
is not something the corpus directly incentivises, and our aggregates (mean, standard deviation,
count) may simply be too coarse to capture it. Second, GRDM's negative lift is best read alongside
a property of its shipped corpus rather than its design: every string and datetime column was
converted to a float and never converted back, so its databases carry no text or categorical
columns at all (0 against RelDiff's 234 over the same references). The RT's text encoder and
MiniLM pathway therefore received no signal in that arm, which is a confound worth stating
plainly: the GRDM arm differs from the others in more than its generator.

\paragraph{Diversity is not the mechanism.}
A natural alternative is that RelDiff simply produces more diverse data. Measuring feature, table
and cross-table diversity directly from the generated databases --- never through any of the four
models, so the measurement cannot be an artifact of what it explains --- RelDiff is indeed most
diverse on all three axes (1.10 / 0.57 / 1.04). But the ranking of the others does not track
mechanism formation: PluRel exceeds GRDM on cross-table diversity (0.85 vs.\ 0.71) yet learns no
usable relational computation (Table~\ref{tab:e7}), while GRDM does on in-distribution schemas.
Diversity describes a generator's marginal output; predictive necessity describes whether values
are \emph{coupled} to structure, and only the latter separates the four models.
